\section{Chỗ vỡ trước không phải phép tính}
\label{sec:ch08-cho-vo}

\index{bộ nhớ!ma trận điểm số}%
Mục trước để lại một câu hỏi: chỗ attention hết rẻ nằm trong cùng một khoảng
\(n\) ở hai bề rộng cách nhau bốn lần, nên nó không phải hàm của \(d\). Mục này
đi tìm nó là hàm của cái gì, và câu trả lời không nằm trong bảng 1.

Bắt đầu từ một đại lượng mà bảng 1 không có cột nào cho: chỗ chứa. Một tầng hồi
quy giữ các kích hoạt hình \((\text{batch}, n, d)\), tăng tuyến tính theo \(n\).
Một tầng attention còn phải dựng \tn{ma trận điểm số}{score matrix}, hình
\((\text{batch}, h, n, n)\):

\begin{equation}
  \text{byte} = \text{batch} \times h \times n^2 \times 4 .
  \label{eq:ch08-byte}
\end{equation}

\index{bộ nhớ!công thức không có d trong đó}%
\textbf{Không có \(d\) trong \eqref{eq:ch08-byte}.} Đó là chỗ đáng dừng lại lâu
hơn một dòng. Mọi số hạng của \eqref{eq:ch08-flop-tang} đều mang \(d\), và ma
trận điểm số thì không: nó vuông theo \(n\), nhân với số head và số câu trong
batch, và bề rộng mô hình biến mất vì \(h \dk = d\) đã bị cộng dồn hết vào rồi.

\begin{measured}{ma trận điểm số của một tầng, batch 8, 8 head, float32}
\begin{minted}{text}
n      bytes               GiB        x a recurrent layer's (b,n,d)
128    4,194,304           0.0039     2.0
512    67,108,864          0.0625     8.0
1024   268,435,456         0.2500     16.0
2048   1,073,741,824       1.0000     32.0
4096   4,294,967,296       4.0000     64.0
8192   17,179,869,184      16.0000    128.0
\end{minted}

Đo bằng \code{experiments/ch08\_flops.py} tại tag \code{ch08}.
\end{measured}

Ở \(n = 2048\) thì đúng một GiB, cho \emph{một} tầng và \emph{một} batch tám
câu, và huấn luyện giữ tất cả chúng sống tới lượt backward. Cột cuối là tỷ số so
với chỗ chứa của một tầng hồi quy cùng cấu hình, và nó nhân đôi mỗi lần \(n\)
nhân đôi.

\subsection*{Một giả thuyết, và cách bác nó}

\index{bộ nhớ!giả thuyết cache}%
Đặt \eqref{eq:ch08-byte} cạnh bảng thời gian ở mục trước thì có một
\tn{giả thuyết}{hypothesis}.
Máy này có 33 MiB cache mức ba. Ở batch 8 và 8 head, \(n = 256\) cho 16 MiB và
\(n = 512\) cho 64 MiB. Chỗ đổi dấu của cột \code{attn/lstm} nằm giữa hai giá
trị ấy, ở \emph{cả hai} bề rộng, và \eqref{eq:ch08-byte} không chứa \(d\) nên nó
cũng không đổi khi \(d\) đổi.

Giả thuyết: cái dốc là chỗ ma trận điểm số thôi không nằm vừa trong cache, và
ngưỡng đo bằng byte chứ không đo bằng \(n\).

Cách kiểm thẳng nhất là bỏ hẳn số tuyệt đối đi và chia thời gian cho số FLOP của
chính phép tính ấy. Một mô hình chi phí chỉ nói về số phép tính thì cột ấy phải
phẳng.

\begin{measured}{giây trên mỗi gigaFLOP, một tầng con attention, \(d = 128\)}
\begin{minted}{text}
n     score MiB   gigaFLOP   seconds     s/GFLOP     vs n=64
64    1.0         0.084      0.000496    0.005912    1.00
128   4.0         0.201      0.000732    0.003638    0.62
192   9.0         0.352      0.001151    0.003267    0.55
256   16.0        0.537      0.002100    0.003912    0.66
320   25.0        0.755      0.004291    0.005684    0.96
384   36.0        1.007      0.009874    0.009809    1.66
448   49.0        1.292      0.015379    0.011905    2.01
512   64.0        1.611      0.022444    0.013935    2.36
640   100.0       2.349      0.035426    0.015082    2.55
\end{minted}

batch 8, 8 head. Cache mức ba của máy là 33 MiB.

Đo bằng \code{experiments/ch08\_clock.py} tại tag \code{ch08}.
\end{measured}

Nó không phẳng. Hàng \(n = 64\) nằm cao vì ở đó phần chi phí cố định vẫn lấn át,
đúng như mọi bảng đồng hồ khác của chương. Từ \(n = 128\) tới \(n = 256\) nó đi
ngang, quanh 0.0033 tới 0.0039 giây mỗi gigaFLOP. Rồi từ \(n = 320\) nó leo, và
leo mạnh: 0.005684, rồi 0.009809, rồi 0.011905, tới 0.015082 ở \(n = 640\). Cùng
một phép tính, cùng một tầng, và mỗi phép tính đắt lên hơn bốn lần so với chỗ
phẳng.

Chỗ gãy nằm giữa 16 MiB và 25 MiB, tức quanh chỗ ma trận điểm số cộng với phần
còn lại của tầng vượt quá 33 MiB.

\subsection*{Phép đối chứng: giữ nguyên \texorpdfstring{\(n\)}{n}, chỉ đổi batch}

\index{bộ nhớ!đối chứng bằng batch}%
Bảng trên một mình chưa kết luận được, và tôi muốn nói rõ vì sao trước khi đi
tiếp. Khi \(n\) tăng thì số byte tăng, nhưng \emph{mọi thứ khác} cũng tăng: số
FLOP, số hàng của softmax, độ dài của mỗi phép quy giản. Một cột leo lên theo
\(n\) có thể là chuyện cache và cũng có thể là năm chuyện khác.

Nên tách hai thứ ra. \eqref{eq:ch08-byte} có ba thừa số, và batch là một trong
ba. Giữ \(n = 256\) cố định, chỉ tăng batch, thì số byte đi qua đúng dải cũ mà
\(n\) không nhúc nhích. Nếu cái dốc đi theo \(n\), nó phải biến mất. Nếu nó đi
theo byte, nó phải xuất hiện lại ở đúng chỗ cũ.

\begin{measured}{cùng số byte, đạt tới bằng batch thay vì bằng \(n\)}
\begin{minted}{text}
batch  score MiB   gigaFLOP   seconds     s/GFLOP
2      4.0         0.134      0.000625    0.004657
4      8.0         0.268      0.001011    0.003766
8      16.0        0.537      0.002777    0.005172
16     32.0        1.074      0.009348    0.008706
32     64.0        2.147      0.023538    0.010961
\end{minted}

\(d = 128\), 8 head, \(n = 256\) ở mọi hàng.

Đo bằng \code{experiments/ch08\_clock.py} tại tag \code{ch08}.
\end{measured}

Nó xuất hiện lại. Từ 4 tới 16 MiB thì cột cuối đi ngang, ba giá trị 0.004657,
0.003766 và 0.005172;
ở 32 MiB nó lên 0.008706; ở 64 MiB lên 0.010961. Cùng một bậc thang, ở cùng một
dải byte, với \(n\) đứng yên suốt cả bảng.

Đặt hai bảng cạnh nhau ở chỗ chúng gặp nhau thì càng rõ. 32 MiB đạt bằng batch
cho 0.008706; 36 MiB đạt bằng \(n\) cho 0.009809. Cùng số byte, cùng bậc tốc độ
mỗi phép tính, mà một bên \(n = 256\) và bên kia \(n = 384\). Còn 64 MiB đạt
bằng batch cho 0.010961 và đúng 64 MiB đạt bằng \(n\) cho 0.013935.

Giả thuyết \enquote{cái dốc đi theo \(n\)} không sống sót. Cái dốc đi theo byte.

\subsection*{Cái này đổi cách đọc bảng 1}

\index{bộ nhớ!bảng 1 không có cột cho chuyện này}%
Ba cột của bảng 1 đều là hàm của \(n\), \(d\), \(k\) và \(r\). Không cột nào là
hàm của kích thước cache, và không cột nào có batch trong đó. Nên bảng ấy không
dự đoán được chỗ máy này chậm lại, không phải vì nó sai mà vì nó không nói về
đại lượng ấy.

Nói cho gọn: \(O(n^2 \cdot d)\) đếm phép tính, và trên phần cứng thật, cái đắt ở
attention không phải phép tính mà là chỗ chứa và việc đi lại giữa các mức bộ
nhớ. Một phép nhân ma trận lớn có nhiều số học trên mỗi byte nó đọc; một lượt
softmax trên ma trận \(n \times n\) thì đọc và ghi đúng chừng ấy byte mà gần như
không có số học nào để bù. Khi ma trận ấy còn nằm trong cache thì chuyện đó
không thấy được. Khi nó không nằm vừa nữa thì mỗi phép tính phải chờ bộ nhớ.

\index{bộ nhớ!hạt giống của FlashAttention}%
Đó cũng là lý do tôi để mục này ở đây thay vì để chương~\ref{ch:chi-phi-bac-hai}
mở đầu bằng nó. FlashAttention không xấp xỉ attention và không giảm số FLOP:
kết quả của nó bằng đúng kết quả của \eqref{eq:ch07-attention}. Cái nó đổi là
không bao giờ dựng ma trận \(n \times n\) ra bộ nhớ, mà chia thành từng khối vừa
cache rồi cộng dồn. Đọc một mình thì đó là một thủ thuật cài đặt. Đọc cạnh hai
bảng trên thì nó là cách sửa đúng vào đại lượng đắt nhất, và bảng 1 không có cột
nào để thấy nó.

Một chú ý cuối về cách tôi đo chỗ chứa, để nếu bạn đo lại thì không mất thời
gian như tôi đã mất: hai cách hiển nhiên đều sai. \code{tracemalloc} báo tám trăm byte ở mọi \(n\) từ
128 tới 4096, phẳng và sai hẳn, vì nó chỉ theo dõi heap của Python còn tensor
của thư viện nằm ngoài đó. Đo bằng chỗ chứa mà tiến trình đang giữ thì đúng bậc
độ lớn nhưng không sạch: mức nền tăng dần qua từng vòng lặp dù đã dọn, vì bộ cấp
phát giữ lại khối đã trả thay vì đưa về hệ điều hành. Nên bảng đầu mục in
\eqref{eq:ch08-byte} chứ không in một phép đo, và bảng thứ hai với thứ ba đo
\emph{thời gian} chứ không đo byte. Số byte thì tính, thời gian thì đo, và
chuyện chúng khớp nhau là kết quả.
